%%%%%%%%%%%%%%%%
%%% deut 7 %%%
%%%%%%%%%%%%%%%%

\Chapter{7}

\Verse{1}
{
	\hDtrA{כִּי יְבִיאֲךָ יְהֹוָה אֱלֹהֶיךָ אֶל הָאָרֶץ אֲשֶׁר אַתָּה בָא שָׁמָּה לְרִשְׁתָּהּ וְנָשַׁל גּוֹיִם רַבִּים מִפָּנֶיךָ הַחִתִּי וְהַגִּרְגָּשִׁי וְהָאֱמֹרִי וְהַכְּנַעֲנִי וְהַפְּרִזִּי וְהַחִוִּי וְהַיְבוּסִי שִׁבְעָה גוֹיִם רַבִּים וַעֲצוּמִים מִמֶּךָּ׃}
}
{
	\eDtrA{
		“When YHWH, your God, will bring you to the land to
		which you’re coming to take possession of it and will
		eject numerous nations from in front of you—the Hittite
		and the Girgashite and the Amorite and the Canaanite and
		the Perizzite and the Hivite and the Jebusite—seven
		nations more numerous and powerful than you,
	}
}
{
%	\footnote{
%		more numerous and powerful than you. Ironically, these are the words
%		that the Pharaoh who enslaved Israel used to describe the Israelites:
%		“Here, the people of the children of Israel is more numerous and
%		powerful than we” (Exod 1:9). But the situation is reversed. The Pharaoh
%		tells his people that they should fear Israel on account of its great
%		size. But God now tells the Israelites not to fear the Canaanite peoples
%		on account of their size.
%	}
}

\Verse{2}
{
	\hDtrA{וּנְתָנָם יְהֹוָה אֱלֹהֶיךָ לְפָנֶיךָ וְהִכִּיתָם הַחֲרֵם תַּחֲרִים אֹתָם לֹא תִכְרֹת לָהֶם בְּרִית וְלֹא תְחנֵּם׃}
}
{
	\eDtrA{
		and YHWH, your God, will put them in front of you, and
		you’ll strike them: you shall completely destroy them! You
		shall not make a covenant with them, and you shall not
		show grace to them.
	}
}
{
%	\footnote{
%		completely destroy. See the comment on Deut 2:34. Many people have been
%		troubled by the idea of commanding the annihilation of the Canaanite
%		residents of the land. The archaeological evidence is that such a
%		destruction never took place. This passage in the Torah was written long
%		after the period of the Israelites’ settlement in the land, and so it is
%		ironic that the author of this text conceived of a degree of violence
%		that appears never in fact to have happened, and then people are
%		troubled by this degree of violence in Israel’s history.
%	}
}

\Verse{3}
{
	\hDtrA{וְלֹא תִתְחַתֵּן בָּם בִּתְּךָ לֹא תִתֵּן לִבְנוֹ וּבִתּוֹ לֹא תִקַּח לִבְנֶךָ׃}
}
{
	\eDtrA{
		And you shall not marry with them. You shall not give your
		daughter to his son, and you shall not take his daughter
		for your son,
	}
}
{
%	\footnote{
%		you shall not marry with them. Moses tells the Israelites that they are
%		no longer free to marry persons from the seven peoples of Canaan. Their
%		ancestors, like Judah (Gen 38:2), could marry Canaanite women because,
%		biblically, a woman takes her husband’s religion. Thus Abraham, Isaac,
%		Jacob, Joseph, Moses, and Samson all marry women who come from other
%		communities but who are understood to take on their husbands’ religion.
%		But Canaanite women are no longer included. There is certainly the
%		possibility that wives will not completely turn to their husbands’
%		religion. This is what happens in the case of Solomon (1 Kings 11) and
%		in the case of Ahab and Jezebel (1 Kings 16:30–33). And this possibility
%		is given as the reason for the ban on marriage to Canaanites, as stated
%		in the next verse (Deut 7:4). This command is curious because in the
%		preceding verse Moses has already said that the Canaanites are to be
%		completely destroyed. The ban on marrying them appears to come from the
%		recognition that the Israelites may fail to eliminate the Canaanites
%		from the land or that it may be a very long process—which is in fact
%		what happens (see Judg 2:1–3).
%	}
}

\Verse{4}
{
	\hDtrA{כִּי יָסִיר אֶת בִּנְךָ מֵאַחֲרַי וְעָבְדוּ אֱלֹהִים אֲחֵרִים וְחָרָה אַף יְהֹוָה בָּכֶם וְהִשְׁמִידְךָ מַהֵר׃}
}
{
	\eDtrA{
		because he’ll turn your son from after me, and they’ll
		serve other gods, and YHWH’s anger will flare at you, and
		He’ll destroy you quickly.
	}
}
{
}

\Verse{5}
{
	\hDtrA{כִּי אִם כֹּה תַעֲשׂוּ לָהֶם מִזְבְּחֹתֵיהֶם תִּתֹּצוּ וּמַצֵּבֹתָם תְּשַׁבֵּרוּ וַאֲשֵׁירֵהֶם תְּגַדֵּעוּן וּפְסִילֵיהֶם תִּשְׂרְפוּן בָּאֵשׁ׃}
}
{
	\eDtrA{
		But this is what you shall do to them: you shall demolish
		their altars and shatter their pillars and cut down their
		Asherahs and burn their statues in fire.
	}
}
{
%	\footnote{
%		Asherahs. See the comment on Deut 16:21.
%	}
}

\Verse{6}
{
	\hDtrA{כִּי עַם קָדוֹשׁ אַתָּה לַיהֹוָה אֱלֹהֶיךָ בְּךָ בָּחַר יְהֹוָה אֱלֹהֶיךָ לִהְיוֹת לוֹ לְעַם סְגֻלָּה מִכֹּל הָעַמִּים אֲשֶׁר עַל פְּנֵי הָאֲדָמָה׃}
}
{
	\eDtrA{
		Because you are a holy people to YHWH, your God. YHWH,
		your God, chose you to become a treasured people to Him
		out of all the peoples who are on the face of the earth.
	}
}
{
%	\footnote{
%		chose you to become a treasured people. As we see in an earlier
%		occurrence of the term “treasured” for Israel, it applies only if they
%		listen and fulfill their covenant (see the comment on Exod 19:5). The
%		same applies to the term “chosen.” It does not refer to any superiority
%		of Israel. The reference to being chosen here is immediately followed by
%		the statement that it is not because of size. And a little farther on,
%		Moses will declare in vivid terms that it is not because of any special
%		merit of Israel’s (Deuteronomy 9). It is because of (a) YHWH’s affection
%		for Israel and (b) YHWH’s oath to the patriarchs. Its practical
%		consequence is that YHWH and Israel have a covenant. That is, being
%		chosen has many sides: it is inspiring, uplifting, a burden, a
%		challenge, a threat, an opportunity.
%	}
}

\Verse{7}
{
	\hDtrA{לֹא מֵרֻבְּכֶם מִכּל הָעַמִּים חָשַׁק יְהֹוָה בָּכֶם וַיִּבְחַר בָּכֶם כִּי אַתֶּם הַמְעַט מִכּל הָעַמִּים׃}
}
{
	\eDtrA{
		It wasn’t because of your being more numerous than all the
		peoples that YHWH was attracted to you so that He chose
		you, because you’re the smallest of all the peoples.
	}
}
{
%	\footnote{
%		attracted. This is the same word that is used to describe Shechem’s
%		longing for Dinah (Gen 34:8) and a soldier’s attraction to a female
%		captive (Deut 21:11). It conveys a feeling of attachment that one soul
%		has to another.
%	}
}

\Verse{8}
{
	\hDtrA{כִּי מֵאַהֲבַת יְהֹוָה אֶתְכֶם וּמִשּׁמְרוֹ אֶת הַשְּׁבֻעָה אֲשֶׁר נִשְׁבַּע לַאֲבֹתֵיכֶם הוֹצִיא יְהֹוָה אֶתְכֶם בְּיָד חֲזָקָה וַיִּפְדְּךָ מִבֵּית עֲבָדִים מִיַּד פַּרְעֹה מֶלֶךְ מִצְרָיִם׃}
}
{
	\eDtrA{
		But because of YHWH’s loving you and because of His
		keeping the oath that He swore to your fathers, YHWH
		brought you out with a strong hand and redeemed you from a
		house of slaves, from the hand of Pharaoh, king of Egypt.
	}
}
{
}

\Verse{9}
{
	\hDtrA{וְיָדַעְתָּ כִּי יְהֹוָה אֱלֹהֶיךָ הוּא הָאֱלֹהִים הָאֵל הַנֶּאֱמָן שֹׁמֵר הַבְּרִית וְהַחֶסֶד לְאֹהֲבָיו וּלְשֹׁמְרֵי מִצְוֺתָו לְאֶלֶף דּוֹר׃}
}
{
	\eDtrA{
		Therefore you shall know that YHWH, your God, He is God,
		the faithful God, keeping the covenant and kindness for
		those who love Him and who observe His commandments to the
		thousandth generation
	}
}
{
}

\Verse{10}
{
	\hDtrA{וּמְשַׁלֵּם לְשֹׂנְאָיו אֶל פָּנָיו לְהַאֲבִידוֹ לֹא יְאַחֵר לְשֹׂנְאוֹ אֶל פָּנָיו יְשַׁלֶּם לוֹ׃}
}
{
	\eDtrA{
		and paying back to those who hate Him to their faces to
		destroy them. He won’t delay toward one who hates Him.
		He’ll pay him back to his face.
	}
}
{
}

\Verse{11}
{
	\hDtrA{וְשָׁמַרְתָּ אֶת הַמִּצְוָה וְאֶת הַחֻקִּים וְאֶת הַמִּשְׁפָּטִים אֲשֶׁר אָנֹכִי מְצַוְּךָ הַיּוֹם לַעֲשׂוֹתָם׃}
}
{
	\eDtrA{
		So you shall observe the commandment and the laws and the
		judgments that I command you today, to do them.
	}
}
{
}

\Verse{12}
{
	\hDtrA{וְהָיָה עֵקֶב תִּשְׁמְעוּן אֵת הַמִּשְׁפָּטִים הָאֵלֶּה וּשְׁמַרְתֶּם וַעֲשִׂיתֶם אֹתָם וְשָׁמַר יְהֹוָה אֱלֹהֶיךָ לְךָ אֶת הַבְּרִית וְאֶת הַחֶסֶד אֲשֶׁר נִשְׁבַּע לַאֲבֹתֶיךָ׃}
}
{
	\eDtrA{
		BECAUSE “And it will be because you’ll listen to these
		judgments and observe and do them that YHWH, your God,
		will keep the covenant and kindness for you that he swore
		to your fathers.
	}
}
{
}

\Verse{13}
{
	\hDtrA{וַאֲהֵבְךָ וּבֵרַכְךָ וְהִרְבֶּךָ וּבֵרַךְ פְּרִי בִטְנְךָ וּפְרִי אַדְמָתֶךָ דְּגָנְךָ וְתִירֹשְׁךָ וְיִצְהָרֶךָ שְׁגַר אֲלָפֶיךָ וְעַשְׁתְּרֹת צֹאנֶךָ עַל הָאֲדָמָה אֲשֶׁר נִשְׁבַּע לַאֲבֹתֶיךָ לָתֶת לָךְ׃}
}
{
	\eDtrA{
		And He’ll love you and bless you and multiply you and
		bless the fruit of your womb and the fruit of your land:
		your grain and your wine and your oil, your cattle’s
		offspring and your flock’s young, on the land that He
		swore to your fathers to give to you.
	}
}
{
%	\footnote{
%		your womb. The word “your” in this phrase is masculine, and so some have
%		interpreted it as meaning that the woman’s womb was thought to belong to
%		her husband. But that is questionable because the entire list of
%		blessings in these verses (13–24) is formulated in the masculine
%		singular, so the reference to the womb was just understood as taking a
%		masculine possessive pronoun like the rest of the list. No conclusions
%		should be drawn from this technical point of grammar regarding ancient
%		Israelite conceptions of male-female ownership of the womb.
%	}
}

\Verse{14}
{
	\hDtrA{בָּרוּךְ תִּהְיֶה מִכּל הָעַמִּים לֹא יִהְיֶה בְךָ עָקָר וַעֲקָרָה וּבִבְהֶמְתֶּךָ׃}
}
{
	\eDtrA{
		You’ll be more blessed than all the peoples. There won’t
		be an infertile male or female among you or among your
		animals.
	}
}
{
%	\footnote{
%		blessed. This word in the Torah generally refers to being well off. Here
%		it is explicitly connected with fertility and good health. So when it
%		says that Israel will be more blessed than other peoples if they keep
%		their covenant, it does not mean that Israel will have some special
%		status, but rather it means that God will bless them with well-being.
%	}
}

\Verse{15}
{
	\hDtrA{וְהֵסִיר יְהֹוָה מִמְּךָ כּל חֹלִי וְכל מַדְוֵי מִצְרַיִם הָרָעִים אֲשֶׁר יָדַעְתָּ לֹא יְשִׂימָם בָּךְ וּנְתָנָם בְּכל שֹׂנְאֶיךָ׃}
}
{
	\eDtrA{
		And YHWH will turn away from you every illness. And all
		the bad diseases of Egypt that you knew: He won’t set them
		among you but will put them among all who hate you.
	}
}
{
}

\Verse{16}
{
	\hDtrA{וְאָכַלְתָּ אֶת כּל הָעַמִּים אֲשֶׁר יְהֹוָה אֱלֹהֶיךָ נֹתֵן לָךְ לֹא תָחוֹס עֵינְךָ עֲלֵיהֶם וְלֹא תַעֲבֹד אֶת אֱלֹהֵיהֶם כִּי מוֹקֵשׁ הוּא לָךְ׃}
}
{
	\eDtrA{
		And you shall eat up all the peoples that YHWH, your God,
		is giving you. Your eye shall not pity them. And you shall
		not serve their gods, because that’s a trap for you.
	}
}
{
}

\Verse{17}
{
	\hDtrA{כִּי תֹאמַר בִּלְבָבְךָ רַבִּים הַגּוֹיִם הָאֵלֶּה מִמֶּנִּי אֵיכָה אוּכַל לְהוֹרִישָׁם׃}
}
{
	\eDtrA{
		“If you say in your heart, ‘These nations are more
		numerous than I; how shall I be able to dispossess them?’
	}
}
{
}

\Verse{18}
{
	\hDtrA{לֹא תִירָא מֵהֶם זָכֹר תִּזְכֹּר אֵת אֲשֶׁר עָשָׂה יְהֹוָה אֱלֹהֶיךָ לְפַרְעֹה וּלְכל מִצְרָיִם׃}
}
{
	\eDtrA{
		you shall not fear them. You shall remember what YHWH,
		your God, did to Pharaoh and to all Egypt:
	}
}
{
}

\Verse{19}
{
	\hDtrA{הַמַּסֹּת הַגְּדֹלֹת אֲשֶׁר רָאוּ עֵינֶיךָ וְהָאֹתֹת וְהַמֹּפְתִים וְהַיָּד הַחֲזָקָה וְהַזְּרֹעַ הַנְּטוּיָה אֲשֶׁר הוֹצִאֲךָ יְהֹוָה אֱלֹהֶיךָ כֵּן יַעֲשֶׂה יְהֹוָה אֱלֹהֶיךָ לְכל הָעַמִּים אֲשֶׁר אַתָּה יָרֵא מִפְּנֵיהֶם׃}
}
{
	\eDtrA{
		the big tests that your eyes saw and the signs and the
		wonders and the strong hand and the outstretched arm by
		which YHWH, your God, brought you out. So YHWH, your God,
		will do to the peoples before whom you fear!
	}
}
{
}

\Verse{20}
{
	\hDtrA{וְגַם אֶת הַצִּרְעָה יְשַׁלַּח יְהֹוָה אֱלֹהֶיךָ בָּם עַד אֲבֹד הַנִּשְׁאָרִים וְהַנִּסְתָּרִים מִפָּנֶיךָ׃}
}
{
	\eDtrA{
		And YHWH, your God, will send the hornet against them as
		well until those who remain and those who are hidden
		perish from in front of you.
	}
}
{
}

\Verse{21}
{
	\hDtrA{לֹא תַעֲרֹץ מִפְּנֵיהֶם כִּי יְהֹוָה אֱלֹהֶיךָ בְּקִרְבֶּךָ אֵל גָּדוֹל וְנוֹרָא׃}
}
{
	\eDtrA{
		Don’t be scared in front of them, because YHWH, your God,
		is among you, a great and fearful God.
	}
}
{
}

\Verse{22}
{
	\hDtrA{וְנָשַׁל יְהֹוָה אֱלֹהֶיךָ אֶת הַגּוֹיִם הָאֵל מִפָּנֶיךָ מְעַט מְעָט לֹא תוּכַל כַּלֹּתָם מַהֵר פֶּן תִּרְבֶּה עָלֶיךָ חַיַּת הַשָּׂדֶה׃}
}
{
	\eDtrA{
		And YHWH, your God, will eject these nations from in front
		of you little by little. You won’t be able to finish them
		quickly, or else the animal of the field will become many
		at you.
	}
}
{
}

\Verse{23}
{
	\hDtrA{וּנְתָנָם יְהֹוָה אֱלֹהֶיךָ לְפָנֶיךָ וְהָמָם מְהוּמָה גְדֹלָה עַד הִשָּׁמְדָם׃}
}
{
	\eDtrA{
		And YHWH, your God, will put them in front of you and put
		them into a big tumult until they’re destroyed.
	}
}
{
%	\footnote{
%		a big tumult. This is the term for what God does to the Egyptians at the
%		Red Sea. A few verses earlier Moses says that God will do to the
%		Canaanites what was done to the Egyptians, and so he now describes what
%		will happen to the Canaanites in the language of the defeat of the
%		Egyptian army.
%	}
}

\Verse{24}
{
	\hDtrA{וְנָתַן מַלְכֵיהֶם בְּיָדֶךָ וְהַאֲבַדְתָּ אֶת שְׁמָם מִתַּחַת הַשָּׁמָיִם לֹא יִתְיַצֵּב אִישׁ בְּפָנֶיךָ עַד הִשְׁמִדְךָ אֹתָם׃}
}
{
	\eDtrA{
		And He’ll put their kings in your hand, and you’ll destroy
		their name from under the skies. Not a man will stand in
		front of you, until you’ve destroyed them.
	}
}
{
}

\Verse{25}
{
	\hDtrA{פְּסִילֵי אֱלֹהֵיהֶם תִּשְׂרְפוּן בָּאֵשׁ לֹא תַחְמֹד כֶּסֶף וְזָהָב עֲלֵיהֶם וְלָקַחְתָּ לָךְ פֶּן תִּוָּקֵשׁ בּוֹ כִּי תוֹעֲבַת יְהֹוָה אֱלֹהֶיךָ הוּא׃}
}
{
	\eDtrA{
		You shall burn the statues of their gods in fire. You
		shall not covet the silver and gold on them and take them
		for yourself, in case you’ll be trapped through it,
		because that is an offensive thing to YHWH, your God,
	}
}
{
}

\Verse{26}
{
	\hDtrA{וְלֹא תָבִיא תוֹעֵבָה אֶל בֵּיתֶךָ וְהָיִיתָ חֵרֶם כָּמֹהוּ שַׁקֵּץ תְּשַׁקְּצֶנּוּ וְתַעֵב תְּתַעֲבֶנּוּ כִּי חֵרֶם הוּא׃}
}
{
	\eDtrA{
		and you shall not bring an offensive thing to your house,
		so that you’ll be a thing to be completely destroyed like
		it. You shall detest it and find it abhorrent, because it
		is a thing to be completely destroyed.
	}
}
{
}
